%-------------------------------------------------------------------------------
%	SECTION TITLE
%-------------------------------------------------------------------------------
\cvsection{Research Experience}

\begin{cventries}


\cventrylongwithdot
    {Research Scientist Intern, FAIR}{Meta}{May 2024 - Aug 2024}{Socially Intelligent Robots under Embodied AI at Fundamental AI Research (FAIR) Labs.}
    {
      \begin{cvitems}
        \item Formulated and developed a method to address personalized adaptation of robot behavior from sparse dialog.
      \end{cvitems}
    }
    
\cventrylongwithdot
    {Research Intern, Mosaic}{Allen Institute for AI (AI2)}{May 2023 - Aug 2023}{Mosaic at AI2, led by Prof. Yejin Choi, conducts research on Commonsense Reasoning in AI systems.}
    {
      \begin{cvitems}
        \item Formulated the problem of goal inference from visual observations of partially complete tasks. 
        \item Developing a method involving commonsense reasoning with language models and visual grounding models for video understanding.
      \end{cvitems}
    }
%-------------------------------------------------------------------------------
%	CONTENT
%-------------------------------------------------------------------------------

% \researchsubsection{Graduate Research Assistant}{Robot Autonomy and Interactive Learning Lab, Georgia Institute of Technology}{Advisor: Prof. Sonia Chernova}

% \begin{cventries}

% \cventrybasic{Proactive Robot Assistance}{Aug 2021 - Present}
% {
% \begin{cvitems}
%     \item Conducting research on Longitudinal Proactive Assistance, towards equipping robots with proactive behaviors rooted in a comprehensive understanding of the daily routines and preferences of their users.
% \end{cvitems}
% }


\end{cventries}

\researchsubsection{The Laboratory for Progress, University of Michigan}{Graduate Researcher}{Advisor: Prof. Chad Jenkins}

\begin{cventries}

\cventrybasic
    {Planning over affordance wayfields}{Jan 2019 - Apr 2019}
    {
      \begin{cvitems}
      \item Explored null-space subsumption architecture to create planners over potential-based affordance representations achieve constrained manipulation tasks. Demonstrated pick and place on a coffee mug without spilling its contents, using OpenRAVE to simulate and visualize.
    \end{cvitems}
    }
    
\cventrybasic
    {Potential Field guided RRTs for motion planning}{Jan 2019 - Apr 2019}
    {
      \begin{cvitems}
      \item Created a hybrid motion planner based on RRT and potential field based gradient descent and achieved greater obstacle-clearance and faster convergence compared to a vanilla RRT.
    \end{cvitems}
    }
    
\cventrybasic
    {Manipulation over Mobile Robot Platform}{Jan 2018 - Apr 2018}
    {
      \begin{cvitems}
      \item Implemented capability for a Fetch robot to open doors using MoveIt package on ROS platform, and integrated it into a mobile manipulation pipeline based on a behavior tree structure.
      \end{cvitems}
    }
\end{cventries}
